\chapter{存储器管理}

\section{存储器管理基本概念}
\concept{管理对象与目标}
\begin{points}
  \item 存储管理的对象是主存储器，即内存。
  \item 主存通常分为系统区和用户区：系统区存放内核程序与数据结构，用户区存放用户进程代码和数据。
  \item 目标：为用户提供方便、安全、充分大的存储空间，支持大型应用和系统程序运行。
\end{points}
\plain{这一节先抓总目标：既要让进程“觉得内存够用又好用”，又要保证不同进程彼此隔离，还要让物理内存尽量被充分利用。}
\exam{简答题常直接问“存储器管理的目标是什么”，按“方便、安全、充分大”三个关键词展开即可。}


\concept{存储器管理四个功能}
\begin{points}
  \item 分配与回收：为进程分配内存，进程结束或换出时回收。
  \item 抽象与映射：让进程看到连续的逻辑地址空间，并把逻辑地址转换成物理地址。
  \item 保护与共享：防止进程非法访问他人空间，同时允许受控共享。
  \item 存储扩充：在逻辑上提供比实际物理内存更大的存储空间。
\end{points}
\plain{这四项功能基本对应后面整章内容：连续/非连续分配解决“怎么分”，地址变换解决“怎么映射”，保护共享解决“怎么不乱”，虚拟存储解决“怎么扩充”。}
\exam{选择题喜欢把某个机制和某个功能错配，比如把分页说成只负责保护不负责映射，要会对应。}


\concept{逻辑地址与物理地址}
\begin{points}
  \item 逻辑地址是 CPU 运行的进程所看到的地址，也叫虚地址或相对地址。
  \item 物理地址是内存中的实际地址，也叫实地址或绝对地址。
  \item 地址变换/地址映射/重定位是把逻辑地址转换为物理地址的过程。
  \item MMU 是常见的硬件支持，负责地址转换和合法性检查。
\end{points}
\plain{进程写代码时面对的是逻辑地址；真正访问内存前，硬件还要把它翻成物理地址。MMU 不只做转换，还会顺手检查“这个地址是不是该进程能访问的”。}
\exam{这部分既会单独考概念，也会作为后面分页、分段、请求分页地址题的前置基础。}


\concept{地址绑定时机}
\begin{points}
  \item 编译时绑定：若装入地址已知，编译器可直接生成绝对地址；位置改变则需重新编译。
  \item 装入时绑定：编译生成可重定位代码，装入时确定物理地址。
  \item 运行时绑定：进程运行中仍可移动，地址在执行时由硬件转换，现代分页/分段常用。
  \item 从源程序到可执行映像通常要经历编译、汇编、链接、装载几个阶段，地址在不同阶段的表示形式不同。
\end{points}
\plain{记忆关键不是死背三个名字，而是分清“什么时候最终确定物理地址”：编译时最早、装入时其次、运行时最灵活。}
\exam{常考“哪种方式需要重新编译”“哪种方式需要硬件支持”“哪种方式允许程序运行时移动”。}


\concept{静态重定位与动态重定位}
\begin{points}
  \item 静态重定位在程序装入内存时一次性完成地址修改，运行中不能随意移动。
  \item 动态重定位在程序运行时进行，每次访问内存时由重定位寄存器或 MMU 转换。
  \item 动态重定位支持进程换入换出、移动和虚拟存储，但需要硬件支持。
  \item 动态重定位的基本形式是：物理地址 $=$ 基址寄存器值 $+$ 逻辑地址。
\end{points}
\plain{静态重定位是“装入时一次改完地址”；动态重定位是“运行时每访问一次都现算一次地址”，所以后者才能支持进程搬家。}
\exam{简答题要能写出两者在“变换时机、是否可移动、是否需硬件支持”上的区别。}


\concept{绝对装入、静态重定位装入与动态重定位装入}
\begin{points}
  \item 绝对装入对应编译时绑定，目标代码中已经是确定物理地址；若装入位置变化，通常必须重新编译。
  \item 静态重定位装入对应装入时绑定，装入程序一次性把逻辑地址改成物理地址，装入后地址固定不变。
  \item 动态重定位装入对应运行时绑定，装入时不改逻辑地址，执行时靠重定位寄存器或 MMU 完成转换。
  \item 三者灵活性依次增强，但对硬件支持的要求也逐步提高。
\end{points}
\plain{这三个“装入方式”本质上就是前三种地址绑定时机在装入阶段的具体落地形式，题目换叫法也要能对应上。}
\exam{选择/简答常混着考“编译时绑定=绝对装入”“装入时绑定=静态重定位装入”“执行时绑定=动态重定位装入”。}


\concept{程序链接方式}
\begin{points}
  \item 静态链接：在程序运行前，把各目标模块和所需库函数链接成一个完整装入模块。
  \item 装入时动态链接：装入内存时再完成部分链接工作，便于多个程序共享同一组库模块。
  \item 运行时动态链接：程序执行过程中按需装入并链接目标模块或共享库，灵活性最高。
  \item 动态链接有利于节省内存和便于库升级，但运行时需要额外的链接支持机制。
\end{points}
\plain{链接解决的是“多个目标模块怎么拼成一个能运行的整体”，重定位解决的是“这个整体放到哪里、地址怎么变”。两者别混。}
\exam{简答题里容易把“动态链接”和“动态重定位”混写，前者处理模块装配，后者处理地址转换。}


\section{连续分配管理}
\concept{单一连续分配}
\begin{points}
  \item 内存用户区一次只装入一道用户程序，用户程序独占用户区。
  \item 实现简单，适合早期单道系统。
  \item 缺点是内存利用率低，不能支持多道程序并发。
\end{points}
\plain{这是最简单的连续分配方式，优点只有“简单”，代价是用户区大部分时间容易空着，吞吐量很低。}
\exam{多见于概念对比题，考它为什么只适合早期单道系统。}


\concept{固定分区分配}
\begin{points}
  \item 系统启动时把内存划分为若干固定大小分区，每个分区装入一个进程。
  \item 分区大小可相等，也可不等。
  \item 优点：实现简单，支持多道程序。
  \item 缺点：分区内部未用空间形成内部碎片；分区数限制多道程序道数。
\end{points}
\plain{固定分区把“分配决策”提前到开机时完成，运行时管理容易，但进程大小和分区大小很难总是正好匹配。}
\exam{常考它为什么产生内部碎片，以及为什么分区数会限制并发进程数。}


\concept{动态分区分配}
\begin{points}
  \item 不预先固定分区，进程装入时根据需要动态划分连续内存区。
  \item 优点：比固定分区灵活，内部碎片少。
  \item 缺点：随着分配和回收，会产生外部碎片，需要紧凑/拼接。
\end{points}
\plain{动态分区是按需切块，所以比固定分区更“省料”，但空闲区会越来越碎，最后常卡在“总量够、连续块不够”。}
\exam{算法题多从这里出，尤其是 FF/NF/BF/WF 的分配结果和外部碎片判断。}


\concept{内部碎片与外部碎片}
\begin{points}
  \item 内部碎片：已经分配给进程但进程用不完的空间，位于已分配区域内部。
  \item 外部碎片：空闲空间总量足够，但分散在各处，无法形成一个足够大的连续区域。
  \item 固定分区和分页容易有内部碎片；动态分区和分段容易有外部碎片。
  \item 直觉：内部碎片是“房间给你了但你没住满”；外部碎片是“空房很多但都太零散，放不下大团队”。
\end{points}
\plain{判断碎片类型时就看“浪费空间是在已分配块里面，还是散落在空闲块之间”；位置不同，处理办法也不同。}
\exam{选择题最爱反着说，尤其把“分页有外部碎片、动态分区有内部碎片”作为干扰项。}


\concept{动态分区分配算法}
\begin{points}
  \item 首次适应 FF：从低地址开始找到第一个足够大的空闲分区。速度较快，但低地址碎片较多。
  \item 循环首次适应 NF：从上次查找结束位置继续查找，分布较均匀。
  \item 最佳适应 BF：选择能满足要求的最小空闲分区，尽量减少剩余空间，但容易产生很多小碎片。
  \item 最坏适应 WF：选择最大的空闲分区，剩余空间仍较大，但可能破坏大空闲区。
\end{points}
\plain{做题别只背定义，要抓“扫描从哪开始、挑哪块、剩余碎片会变什么样”这三件事。}
\exam{算法题要把每次分配后的空闲分区表写出来，否则很容易在 NF 或 BF 上算错。}


\concept{动态分区回收与合并}
\begin{points}
  \item 进程释放分区后，需要把回收区按地址顺序插回空闲分区表或空闲分区链。
  \item 若回收区上方邻接空闲分区，则与上邻空闲分区合并；若下方邻接空闲分区，则与下邻空闲分区合并。
  \item 若上下两侧都邻接空闲分区，则三者合并成一个更大的连续空闲区。
  \item 若上下都不邻接空闲分区，则为该回收区单独建立新空闲表项。
  \item 因而回收后空闲分区个数可能不变、减少 1，或增加 1，要根据相邻情况判断。
\end{points}
\plain{动态分区回收题本质是在判“回收区左右有没有空闲邻居”，再决定是并到旧块里还是新开一个空闲块。}
\exam{PPT 专门分四种相邻情况讲回收过程，题目常要求更新空闲分区表，并判断回收后空闲分区数怎么变化。}


\concept{伙伴系统}
\begin{points}
  \item 伙伴系统按 $2^k$ 大小管理内存块，申请时分裂大块，释放时若相邻伙伴空闲则合并。
  \item 优点：分裂和合并规则简单，管理效率较高。
  \item 缺点：申请大小不一定正好是 $2^k$，可能产生内部碎片。
\end{points}
\plain{伙伴系统的关键不是“按 2 的幂分配”这句话本身，而是你能不能顺着分裂和回收过程把内存图画出来。}
\exam{老师明确提过会考画图题：给定总内存和若干申请/释放，按分裂与合并顺序画出每一步。}


\concept{伙伴地址公式}
\begin{points}
  \item 若某块大小为 $2^k$，起始地址为 $d$，则其伙伴块大小相同，起始地址与它只相差 $2^k$。
  \item 若以地址 0 为基准编号，可写成：当 $d \bmod 2^{k+1}=0$ 时，伙伴地址为 $d+2^k$；当 $d \bmod 2^{k+1}=2^k$ 时，伙伴地址为 $d-2^k$。
  \item 若整段可分配空间不是从 0 开始，而是从基址 $a$ 开始，则把上式中的 $d$ 改为相对地址 $d-a$ 再判断。
  \item 直觉上就是看当前块处在某个 $2^{k+1}$ 大块的前半段还是后半段：前半段找后半段，后半段找前半段。
\end{points}
\plain{伙伴地址题本质是在做“同一对兄弟块之间的切换”：块大小不变，只把起始地址在同一更大块的前后两半之间来回翻。}
\exam{PPT 单独给了地址公式和二进制地址例题，考试若出“块大小分别为 4 和 16 时伙伴地址是多少”，先确定 $k$，再按前半段/后半段判断即可。}


\section{分页存储管理}
\concept{分页基本思想}
\begin{points}
  \item 将逻辑地址空间划分为大小相等的页 page。
  \item 将物理内存划分为同样大小的页框/页帧 frame。
  \item 以页为单位分配内存，逻辑上连续，物理上可以分散。
  \item 分页取消了作业对物理连续内存的要求，避免外部碎片，但最后一页可能有内部碎片。
\end{points}
\plain{分页的出发点就是把“必须找一整块连续内存”改成“找若干个页框就行”，所以它主要解决的是连续分配下的外部碎片问题。}
\exam{本章重点在分页，至少要会说清楚它为什么能避免外部碎片、为什么又还会留下内部碎片。}


\concept{页表}
\begin{points}
  \item 系统为每个进程建立一张页表，记录逻辑页号到物理块号的映射。
  \item 每个页面对应一个页表项，页表项通常包含物理块号和标志位。
  \item 标志位可能包括存在位、修改位、引用位、保护位等。
  \item 进程切换时，页表基址和长度等信息会被装入页表寄存器。
\end{points}
\plain{页表就是分页系统里的“翻译词典”：逻辑页号是索引，页表项给出物理块号和访问控制信息。}
\exam{除地址计算外，还常考页表是不是系统共享一张，以及页表项里常见标志位的含义。}


\concept{分页地址结构与转换}
\begin{points}
  \item 逻辑地址由页号 $p$ 和页内偏移 $d$ 组成。
  \item 页大小为 $2^n$ 字节时，低 $n$ 位是页内偏移，高位是页号。
  \item 地址转换：用页号查页表得到物理块号 $f$，物理地址为 $f\times$ 页大小 $+d$。
  \item 若页表项无效或越界，则产生地址越界或缺页异常。
\end{points}
\plain{分页地址题就是三步：拆页号和偏移，查页表，拼物理地址。}
\exam{步骤题要把“先比页号和页表长度、越界则中断”也写出来，老师 PPT 专门强调过合法性检查。}

\concept{分页系统中的碎片}
\begin{points}
  \item 分页消除了外部碎片，因为页框可离散分配。
  \item 但最后一页往往装不满，因此仍可能存在内部碎片。
  \item 内部碎片通常不超过一页大小。
\end{points}
\plain{分页把“大块连续放不下”的问题解决了，但最后那一页角落里可能还有点空着。}
\exam{常考分页有没有外部碎片、为什么还有内部碎片。}

\concept{快表 TLB 与有效访问时间}
\begin{points}
  \item 快表 TLB 是页表项的高速缓存，用于减少访问页表的内存次数。
  \item 无快表时，访问一次数据通常需先访问页表，再访问数据，共两次内存访问。
  \item 有快表时，若命中，直接得到物理块号，只需访问数据；若未命中，需要访问页表再访问数据。
  \item 有效访问时间公式常写为：$EAT=\alpha(t_{TLB}+t_M)+(1-\alpha)(t_{TLB}+2t_M)$，其中 $\alpha$ 为命中率。
\end{points}
\plain{TLB 的作用就是把“每次都查内存里的页表”变成“尽量先查一个更快的小缓存”，所以 EAT 题本质上是在算命中和未命中的加权平均。}
\exam{EAT 是老师点名的重点，做题先分清命中路径和未命中路径各访问几次内存，再代公式。}


\concept{页表实现要点}
\begin{points}
  \item 进程创建时初始化页表，并把页表基址和页表长度保存到 PCB 中。
  \item 进程切换时，内核把当前进程的页表基址、长度等信息装入相应寄存器。
  \item 因而页表是“每个进程一张”，不是整个系统只共用一张。
\end{points}
\plain{页表本质上是“这个进程自己的地址翻译字典”，切换进程就得切换这本字典。}
\exam{常考页表属于谁、什么时候切换页表寄存器。}


\concept{多级页表}
\begin{points}
  \item 多级页表把页表本身分页，只为实际用到的页表页分配内存。
  \item 解决页表过大、必须连续存放的问题。
  \item 缺点：地址转换需要多次查表，若无 TLB 支持会增加访问开销。
  \item 做题时根据每级索引位数判断每级表项数，根据页大小和表项大小判断一页能装多少表项。
\end{points}
\plain{多级页表是“用时间换空间”：多查几次表，换来页表本身不必一次性整块常驻内存。}
\exam{计算题常把“页大小、表项大小、逻辑地址位数”混在一起考，要先算一页能容纳多少表项，再反推需要几级。}

\concept{分页存储管理的保护与共享}
\begin{points}
  \item 地址越界保护：将逻辑地址中的页号与页表长度比较，页号越界则产生中断。
  \item 存取控制保护：在页表项中设置保护位，限制页面的读、写、执行等访问方式。
  \item 数据共享：不同进程可把各自页表中的某些页号映射到同一物理页框，实现共享数据页。
  \item 代码共享：若共享代码页中含有逻辑地址，则不同进程通常需把该代码页放在各自逻辑空间中的相同页号处。
\end{points}
\plain{分页下的共享本质是“多张页表指向同一页框”；分页下的保护本质是“先防越界，再靠页表项权限位限制怎么访问”。}
\exam{简答题若问分页如何实现保护与共享，答题结构直接写“越界保护、存取控制保护、数据共享、代码共享条件”最稳。}


\section{分段与段页式存储管理}
\concept{分段基本思想}
\begin{points}
  \item 分段按程序的逻辑结构划分地址空间，如代码段、数据段、栈段、共享库段。
  \item 逻辑地址由段号和段内偏移组成。
  \item 段表项包含段基址、段长和保护信息。
  \item 地址转换时先检查段号和段内偏移是否合法，再由段基址加偏移得到物理地址。
\end{points}
\plain{分段这节最该抓的是“先看逻辑意义再谈地址”：代码段、数据段、栈段天然分开，所以共享和保护都更符合程序员直觉。}
\exam{常考分段地址题时先做越界判断，再做“段基址 + 段内偏移”；若一上来只会相加，通常会丢步骤分。}

\concept{分段保护和共享}
\begin{points}
  \item 地址越界保护包括两层：先比较段号与段表长度，再比较段内位移与该段段长。
  \item 存取控制保护：可按段统一设置读、写、执行权限，因为同一段内信息通常具有相同访问属性。
  \item 分段天然便于共享，因为段本身就是逻辑信息单位。
  \item 共享实现方式：让多个进程的段表项指向同一段在内存中的起始地址即可。
\end{points}
\plain{分段比分页更适合讲“共享”和“权限”，因为它按代码段、数据段、栈段这类逻辑单位划分，管理粒度更贴合程序结构。}
\exam{问“为什么分段比分页更便于实现共享和保护”时，要抓住“段是逻辑单位、可按段统一赋权、多个段表项可共指同一内存段”三点。}


\concept{分页与分段比较}
\begin{longtable}{p{0.20\textwidth}p{0.36\textwidth}p{0.36\textwidth}}
\toprule
比较项 & 分页 & 分段 \\
\midrule
划分依据 & 系统按固定大小机械划分 & 用户程序逻辑结构划分 \\
单位大小 & 页大小固定 & 段大小可变 \\
地址空间 & 一维地址空间 & 二维地址空间：段号+段内偏移 \\
碎片 & 无外部碎片，有内部碎片 & 无内部碎片或较少，有外部碎片 \\
共享保护 & 不如分段自然 & 以逻辑段为单位，便于共享与保护 \\
\bottomrule
\end{longtable}

\concept{段页式管理}
\begin{points}
  \item 段页式先把程序按逻辑分段，再把每个段分页。
  \item 地址结构通常为段号、段内页号、页内偏移。
  \item 每个段有段表项，段表项指向该段对应的页表。
  \item 优点：兼具分段的逻辑保护共享和分页的离散分配；缺点是地址转换更复杂。
\end{points}
\plain{段页式不要背成三个名词堆在一起，它本质是在说：先借“段”保留逻辑结构，再借“页”解决离散分配问题。}
\exam{常考段页式地址结构和查表顺序，答题时最好明确写出“段号 -> 段表 -> 页表 -> 物理块号 -> 页内偏移”。}

\concept{段页式地址题的做法}
\begin{points}
  \item 先从逻辑地址中拆出段号、段内页号、页内偏移。
  \item 用段号查段表，得到该段页表的起始位置并检查段是否合法。
  \item 再用段内页号查该段对应页表，得到物理块号。
  \item 最后由物理块号和页内偏移拼出物理地址。
\end{points}
\plain{段页式就是“先定位哪一段，再定位这一段里的哪一页”，顺序不能反。}
\exam{这类题老师录音特别提过，步骤一定要写完整：拆地址、查段表、查页表、拼地址。}
